\documentclass[10pt,letterpaper]{article}
% Shared package block. Input from the preamble, before palette.tex.
\usepackage[margin=2cm]{geometry}
\usepackage[T1]{fontenc}
\usepackage[utf8]{inputenc}
\usepackage{lmodern}
\usepackage{microtype}
\usepackage{graphicx}
\usepackage{caption}
\usepackage{subcaption}
\usepackage{hyperref}
\usepackage{enumitem}
\usepackage{titlesec}
\usepackage{xcolor}
\usepackage{float}
\usepackage{amssymb}
\graphicspath{{shots/}{figures/}{images/}}

\def\pdfkind{brief}
% Color palette. Requires \pdfkind already defined.
% Kinds: review, release, finance, update, plan, brief.
% Names defined here: accent, onaccent, muted, link, rule.
\makeatletter
\newcommand{\PdfSetPalette}[5]{%
  \definecolor{accent}{HTML}{#1}%
  \definecolor{onaccent}{HTML}{#2}%
  \definecolor{muted}{HTML}{#3}%
  \definecolor{link}{HTML}{#4}%
  \definecolor{rule}{HTML}{#5}%
}
\@namedef{pdfpalette@review}{\PdfSetPalette{F07A1F}{FFFFFF}{5C534C}{1F4E79}{F6D7BE}}
\@namedef{pdfpalette@release}{\PdfSetPalette{C2410C}{FFFFFF}{57534E}{9A3412}{FED7AA}}
\@namedef{pdfpalette@finance}{\PdfSetPalette{0F6E56}{FFFFFF}{3F4F4A}{0F6E56}{C6E4D8}}
\@namedef{pdfpalette@update}{\PdfSetPalette{1D4E89}{FFFFFF}{4A5562}{1D4E89}{D3DDEA}}
\@namedef{pdfpalette@plan}{\PdfSetPalette{3F3F46}{FFFFFF}{52525B}{1E3A5F}{E4E4E7}}
\@namedef{pdfpalette@brief}{\PdfSetPalette{111827}{FFFFFF}{4B5563}{1D4ED8}{E5E7EB}}
\newcommand{\PdfApplyKind}[1]{%
  \@ifundefined{pdfpalette@#1}{%
    \PackageWarning{palette}{Unknown pdfkind '#1', using brief}%
    \@nameuse{pdfpalette@brief}%
  }{%
    \@nameuse{pdfpalette@#1}%
  }%
}
\makeatother
\expandafter\PdfApplyKind\expandafter{\pdfkind}

% Layout macros. Input after palette.tex so accent/onaccent/muted/link/rule exist.
\hypersetup{colorlinks=true,urlcolor=link,linkcolor=link}
\titleformat{\section}{\large\bfseries\color{accent}}{}{0em}{}
\titleformat{\subsection}{\normalsize\bfseries\color{accent}}{}{0em}{}
\captionsetup{font=small,skip=6pt}
\setlist[itemize]{leftmargin=1.2em,itemsep=0.3em}
\setlist[enumerate]{leftmargin=1.6em,itemsep=0.15em,topsep=0.2em}

\newcommand{\colorrule}{\par\noindent{\color{rule}\rule{\linewidth}{2pt}}\par}

% Numbered item: \docitem{title}{meta line}
\newcounter{docitem}
\newcommand{\docitem}[2]{%
  \par\bigskip\refstepcounter{docitem}%
  \noindent{\large\bfseries\color{accent}\thedocitem.}\ {\large\bfseries #1}\par\nopagebreak
  \noindent{\small\color{muted}#2}\par\nopagebreak\smallskip}

% Phone or portrait figure beside the text: \withphone{file}{caption}{text}
\newcommand{\withphone}[3]{%
  \par\noindent
  \begin{minipage}[t]{0.6\linewidth}\vspace{0pt}#3\end{minipage}\hfill
  \begin{minipage}[t]{0.34\linewidth}\vspace{0pt}\centering
    \includegraphics[width=\linewidth,height=8cm,keepaspectratio]{#1}
    \captionof{figure}{#2}
  \end{minipage}\par}

% Still frame beside the text; frame and button open the video.
% \withvideo{frame.png}{caption}{https://...}{text}
\newcommand{\withvideo}[4]{%
  \par\noindent
  \begin{minipage}[t]{0.6\linewidth}\vspace{0pt}#4\end{minipage}\hfill
  \begin{minipage}[t]{0.34\linewidth}\vspace{0pt}\centering
    \href{#3}{\includegraphics[width=\linewidth,height=8cm,keepaspectratio]{#1}}\\[0.3em]
    \href{#3}{\colorbox{accent}{\textcolor{onaccent}{\small\bfseries\ $\blacktriangleright$\ Watch the video\ }}}
    \captionof{figure}{#2}
  \end{minipage}\par}

% Wide figure under the text (dashboards, tables-as-images).
\newcommand{\wide}[2]{%
  \begin{figure}[H]\centering
    \includegraphics[width=\linewidth,keepaspectratio]{#1}
    \caption{#2}
  \end{figure}}


\usepackage{array}
\usepackage{booktabs}
\usepackage{colortbl}
\usepackage{longtable}
\usepackage{listings}
\usepackage{tabularx}
\usepackage{fancyhdr}
\usepackage{lastpage}

\geometry{top=1.6cm,bottom=1.7cm,left=1.75cm,right=1.75cm}
\setlength{\parindent}{0pt}
\setlength{\parskip}{0.35em}
\setlength{\LTpre}{0.2em}
\setlength{\LTpost}{0.6em}
\renewcommand{\arraystretch}{1.16}
\pagestyle{fancy}
\fancyhf{}
\fancyfoot[C]{\small Page \thepage\ of \pageref{LastPage}}
\renewcommand{\headrulewidth}{0pt}

\lstdefinestyle{sqlstyle}{
  language=SQL,
  basicstyle=\ttfamily\scriptsize,
  keywordstyle=\bfseries\color{accent},
  commentstyle=\itshape\color{muted},
  stringstyle=\color{link},
  backgroundcolor=\color{rule!26},
  frame=single,
  rulecolor=\color{rule},
  framerule=0.5pt,
  breaklines=true,
  breakatwhitespace=false,
  columns=fullflexible,
  keepspaces=true,
  showstringspaces=false,
  tabsize=2,
  aboveskip=0.35em,
  belowskip=0.55em
}
\lstset{style=sqlstyle}

\newcommand{\fieldtable}[1]{%
  \begin{tabularx}{\linewidth}{>{\bfseries\columncolor{rule!52}}p{0.27\linewidth}|X}
    \rowcolor{accent}\color{onaccent}Field & \color{onaccent}Details\\
    #1
  \end{tabularx}}

\newcommand{\querymeta}[2]{%
  \begin{tabularx}{\linewidth}{>{\bfseries\columncolor{rule!52}}p{0.22\linewidth}|X}
    \rowcolor{accent}\color{onaccent}Field & \color{onaccent}Details\\
    Category & #1\\
    Purpose & #2\\
  \end{tabularx}}

\newcommand{\queryoutput}[1]{%
  \textbf{Output (paste a screenshot, or a text table for SELECT queries):}\par
  \IfFileExists{images/revised/query-#1-output.png}{%
    \begin{center}
      \includegraphics[width=0.97\linewidth,height=7cm,keepaspectratio]{images/revised/query-#1-output.png}
    \end{center}
  }{%
    \begin{center}
      \fcolorbox{rule}{white}{\parbox[c][5.2cm][c]{0.93\linewidth}{\centering
        \color{muted}\itshape Paste output screenshot here.\\[0.5em]
        Expected file: \texttt{images/revised/query-#1-output.png}}}
    \end{center}
  }}

\newcommand{\querytitle}[1]{\subsection*{Query #1}}

\begin{document}

\begin{center}
  {\Large\bfseries SOUTHEAST UNIVERSITY}\par
  {\large\bfseries CSE384 -- Database Design Lab}\par
  {\large\bfseries\color{accent} PROJECT SQL QUERIES SUBMISSION}\par
\end{center}
\colorrule

\textit{Do not rename, remove or reorder the section headings below. Replace the bracketed placeholders only. Submission is evaluated automatically and then reviewed.}

\section{Section 1: Student Information}

\fieldtable{
  Student Name & [ your full name ]\\
  Student ID & [ your ID ]\\
  Section & [ section ]\\
  Group ID & [ group number, or `Individual' ]\\
  Project Title & DropX -- Parcel Delivery \& Logistics Management System\\
  DBMS Used & MySQL 8.0+\\
}

\section{Section 2: Peer Review and Contribution}

List every member of your group (including yourself) and the percentage of total project work each member contributed. The total must equal 100\%. Every member submits this table independently, so answers are cross-checked.

\begin{tabularx}{\linewidth}{|p{0.22\linewidth}|p{0.16\linewidth}|X|p{0.14\linewidth}|}
  \hline
  \rowcolor{accent}\color{onaccent}\bfseries Member Name & \color{onaccent}\bfseries Student ID & \color{onaccent}\bfseries Work Done (1 line) & \color{onaccent}\bfseries Contribution \%\\
  \hline
  [ name ] & [ ID ] & [ e.g. designed schema, wrote report queries ] & [ \% ]\\ \hline
  [ name ] & [ ID ] & [ ... ] & [ \% ]\\ \hline
  [ name ] & [ ID ] & [ ... ] & [ \% ]\\ \hline
  [ name ] & [ ID ] & [ ... ] & [ \% ]\\ \hline
  \textbf{TOTAL} & & & \textbf{100\%}\\ \hline
\end{tabularx}

\vspace{0.6em}
\textbf{Peer comment (2--3 lines, honest feedback on the teamwork, including any issue you faced):}

\fcolorbox{rule}{white}{\parbox[c][1.4cm][c]{0.96\linewidth}{\color{muted}\itshape [ write here ]}}

\textbf{Self-declaration: I confirm the information above is accurate and the queries below are from our own project.}

\section{Section 3: Database Schema}

Paste the complete CREATE TABLE statements (with all constraints) and sample INSERT statements, so the queries below can be re-run exactly as written.

\textbf{Complete CREATE TABLE statements}
\lstinputlisting{schema-create-tables.sql}

\textbf{Sample INSERT statements}
\lstinputlisting{sample-data.sql}

\section{Section 4: Project Queries}

Use the exact block format below for every query. Queries are numbered continuously and use only the categories allowed by the supplied template.

\textit{Category must be one of: DDL / INSERT / UPDATE / DELETE / SELECT-Filter / ORDER BY / LIKE / IN-BETWEEN / TOP-LIMIT / JOIN / Aggregate-GROUP BY / Subquery / Other.}

\querytitle{1}
\querymeta{Other}{Display the real MySQL definition of the pricing slab table so its primary key, unique keys, foreign key, and CHECK constraint can be verified.}

\textbf{SQL:}
\begin{lstlisting}
SHOW CREATE TABLE pricing_slabs;
\end{lstlisting}
\queryoutput{01}

\clearpage
\querytitle{2}
\querymeta{INSERT}{Record a new rider application submitted through the DropX rider recruitment flow.}

\textbf{SQL:}
\begin{lstlisting}
INSERT INTO rider_applications
    (name, phone, email, district, vehicle_type,
     license_number, experience_years, availability, notes)
VALUES
    ('Imran Kabir', '01600000001',
     'imran.rider@example.com', 'Dhaka', 'MOTORCYCLE',
     'DL-DHK-20002', 2.5, 'Full-time',
     'Available for morning and evening shifts');
\end{lstlisting}
\queryoutput{02}

\clearpage
\querytitle{3}
\querymeta{UPDATE}{Activate a temporary customer after successful OTP verification and store the activation time.}

\textbf{SQL:}
\begin{lstlisting}
UPDATE customers
SET status = 'ACTIVE',
    activated_at = CURRENT_TIMESTAMP
WHERE id = 4
  AND status = 'TEMP';
\end{lstlisting}
\queryoutput{03}

\clearpage
\querytitle{4}
\querymeta{DELETE}{Delete a customer's saved booking draft after the booking is completed or intentionally discarded.}

\textbf{SQL:}
\begin{lstlisting}
DELETE FROM parcel_drafts
WHERE customer_id = 1;
\end{lstlisting}
\queryoutput{04}

\clearpage
\querytitle{5}
\querymeta{SELECT-Filter}{Find the customer account matching the submitted phone number or email during OTP login.}

\textbf{SQL:}
\begin{lstlisting}
SELECT id, name, phone, email, status
FROM customers
WHERE phone = '01900000001'
   OR email = 'ayesha@example.com'
LIMIT 1;
\end{lstlisting}
\queryoutput{05}

\clearpage
\querytitle{6}
\querymeta{ORDER BY}{Show a parcel's tracking history with the newest event first and its hub name when available.}

\textbf{SQL:}
\begin{lstlisting}
SELECT e.event_type, e.description, e.created_at,
       h.name AS hub_name
FROM parcel_events AS e
LEFT JOIN hubs AS h ON h.id = e.hub_id
WHERE e.parcel_id = 1
ORDER BY e.created_at DESC, e.id DESC
LIMIT 100;
\end{lstlisting}
\queryoutput{06}

\clearpage
\querytitle{7}
\querymeta{LIKE}{Find customers whose names contain ``Rahman'' when staff only know part of the customer name.}

\textbf{SQL:}
\begin{lstlisting}
SELECT c.code, c.name, c.phone, c.email, c.status
FROM customers AS c
WHERE c.name LIKE '%Rahman%'
ORDER BY c.name ASC;
\end{lstlisting}
\queryoutput{07}

\clearpage
\querytitle{8}
\querymeta{IN-BETWEEN}{Find parcels in selected active statuses that were created during the submission's sample date range.}

\textbf{SQL:}
\begin{lstlisting}
SELECT p.tracking_number, p.receiver_name, p.status,
       p.delivery_fee, p.created_at
FROM parcels AS p
WHERE p.status IN ('CREATED', 'AT_HUB', 'OUT_FOR_DELIVERY')
  AND p.created_at BETWEEN '2026-10-10 00:00:00'
                       AND '2026-10-14 23:59:59'
ORDER BY p.created_at DESC;
\end{lstlisting}
\queryoutput{08}

\clearpage
\querytitle{9}
\querymeta{TOP-LIMIT}{Display the five newest rider applications for the admin review queue. MySQL uses LIMIT instead of TOP.}

\textbf{SQL:}
\begin{lstlisting}
SELECT id, name, phone, district, vehicle_type,
       status, created_at
FROM rider_applications
ORDER BY created_at DESC, id DESC
LIMIT 5;
\end{lstlisting}
\queryoutput{09}

\clearpage
\querytitle{10}
\querymeta{JOIN}{Combine a parcel with its origin, destination, and current hub to produce the public tracking summary.}

\textbf{SQL:}
\begin{lstlisting}
SELECT p.tracking_number, p.parcel_type,
       p.payment_type, p.status, p.weight, p.cod_amount,
       origin.name AS origin_hub,
       destination.name AS destination_hub,
       current_hub.name AS current_hub
FROM parcels AS p
INNER JOIN hubs AS origin ON origin.id = p.origin_hub_id
INNER JOIN hubs AS destination
    ON destination.id = p.destination_hub_id
LEFT JOIN hubs AS current_hub
    ON current_hub.id = p.current_hub_id
WHERE p.tracking_number = 'DX-20261001-001'
LIMIT 1;
\end{lstlisting}
\queryoutput{10}

\clearpage
\querytitle{11}
\querymeta{Aggregate-GROUP BY}{Count how many parcels are assigned to each selected transfer for the transfer list.}

\textbf{SQL:}
\begin{lstlisting}
SELECT transfer_id, COUNT(*) AS parcel_count
FROM transfer_parcels
WHERE transfer_id IN (1)
GROUP BY transfer_id;
\end{lstlisting}
\queryoutput{11}

\clearpage
\querytitle{12}
\querymeta{Subquery}{Return each parcel with its latest successful delivery time calculated by a correlated subquery.}

\textbf{SQL:}
\begin{lstlisting}
SELECT p.tracking_number, p.status,
       (
           SELECT MAX(d.delivered_at)
           FROM deliveries AS d
           WHERE d.parcel_id = p.id
             AND d.status = 'DELIVERED'
       ) AS delivered_at
FROM parcels AS p
ORDER BY p.tracking_number;
\end{lstlisting}
\queryoutput{12}

\end{document}
